\documentclass[11pt]{article}
\usepackage[a4paper,margin=1in]{geometry}
\usepackage{amsmath,amssymb,amsthm,booktabs,microtype}
\usepackage[hidelinks]{hyperref}
\newtheorem{proposition}{Proposition}
\newcommand{\norm}[1]{\left\lVert #1\right\rVert}
\title{Source-First Growing-Shell Stability of the Native\
Weighted/Positive Budget Gap}
\author{Liang Wang\\
School of Mathematics and Statistics\\
Huazhong University of Science and Technology (HUST), Wuhan, China}
\date{29 August 2026}

\begin{document}
\maketitle

\begin{abstract}
The preceding finite native-profile audit found a robust budget separation
between a signed weighted target and an all-positive control, but its target
labels were available on only 18 inherited rows.  We extend the hostile test
to the 34-row growing/control grid of the prime-shell dynamical model and
recompute the weighted target from the physical output Gram matrix on every
row.  The finite Gram and sign-enumeration identities are exact.  At relative
RMS tolerances $1/4$, $1/2$, and $3/4$, the common-prefix weighted/positive
budget gap exceeds $10$ on all 34 rows; its minima are $85.3204$, $38.2187$,
and $39.2637$.  The common weighted budget exceeds $10^{-5}$ in all 102
row--tolerance cases under each of three source normalizations.  This is a
source-first finite growing-grid certificate.  It does not prove uniform
profile-budget growth, arithmetic $L^2$, or the twin-prime conjecture.
\end{abstract}

\section{Question and finite model}

TPC-301 showed that a finite native obstruction survives a tolerance ladder,
a common-prefix comparison, and three source normalizations \cite{tpc301}.
The unresolved question was whether that obstruction depended on the inherited
18-row target atlas.  Here the target itself is regenerated on the complete
TPC-288 growth/control grid.

Let $S$ be a finite prime shell and let $g_q$ be the exact literal physical
output vector attached to $q\in S$.  Put
\[
 G_{q,r}=\langle g_q,g_r\rangle,\qquad
 R(a)=\frac{a^TGa}{\operatorname{tr}G},\quad a_q\in\{\pm1\}.
\]
The weighted target $a_w$ is the exact minimizer of $R$ after fixing one
global sign; the positive control is $a_+=(1,\ldots,1)$.  The source profiles
are the first 17 literal cutoff profiles
\[
 u_z(t)=\lambda(t)-\sum_{d\le z,\ d\mid t}\mu(d),
 \quad z=3,5,7,11,13,17,19,23,29,31,37,41,43,47,53,59,61.
\]

For the first $k$ profiles, write $V_k=A^TU_k$ and $M_k=U_k^TU_k$.  We use
the relative target budget
\[
 B_{k,\tau}(b)=\min\{c^TM_kc:\norm{V_kc-b}_2\leq\tau\norm b_2\}.
 \tag{1}
\]
The grid has 16 growth-path rows and 18 source-control rows.  Their shell
lists contain 430 explicit targets.  The parent metadata also contains a
1,380-edge count for a different inherited census; we keep the two numbers
separate.

\section{Exact finite identities}

\begin{proposition}[Gram identity and global signs]
For every finite shell, $G$ is positive semidefinite and $R(a)=R(-a)$.
Fixing the first sign to $+1$ therefore represents every global-sign class.
\end{proposition}
\begin{proof}
For any coefficient vector $c$,
\[
 c^TGc=\sum_{q,r}c_qc_r\langle g_q,g_r\rangle
       =\left\|\sum_qc_qg_q\right\|_2^2\geq0.
\]
The quadratic form contains each sign twice, so replacing every sign by its
negative leaves it unchanged.
\end{proof}

\begin{proposition}[Finite enumeration]
After clearing a common positive denominator, a reflected Gray traversal of
the $2^{|S|-1}$ tail bit strings gives the exact minimum of $R$ on the
equal-sign domain.
\end{proposition}
\begin{proof}
The traversal flips one tail coordinate at a time and visits every tail bit
string once.  The incremental quadratic update is an algebraic rearrangement
of the direct integer quadratic form, so no floating-point comparison is used
in selecting the minimum.
\end{proof}

\begin{proposition}[Budget monotonicity and normalization]
Relaxing $\tau$ or enlarging a profile prefix cannot increase (1).  At a
common prefix, a positive target-independent source normalizer cancels from
the weighted/positive budget ratio.
\end{proposition}
\begin{proof}
Both statements follow by inclusion of feasible sets.  For normalization,
divide both numerator and denominator by the same positive scalar.
\end{proof}

\section{Source-first audit}

For every row we rebuild the source weights and all physical output vectors
with rational arithmetic, form $G$, and enumerate $2^{|S|-1}$ sign classes.
The resulting $a_w$ is then inserted into (1), alongside $a_+$.  At each
$\tau\in\{1/4,1/2,3/4\}$ we retain three contexts: the smallest common
feasible prefix, target-specific first feasible prefixes, and the full prefix.
The primary comparison is the common-prefix ratio.  The source normalizers
are $\norm\beta_2^2$, $\operatorname{tr}(M_k)/k$, and $M_k[1,1]$.

The frontier is solved by 60-digit ridge-parameter bisection.  Decimal
intervals are written outward; exact rational labels and ratios are retained
in the canonical certificate.  An independent checker rebuilds all 34 Gram
matrices without importing this producer.

\section{Results}

\begin{table}[t]
\centering\small
\caption{Minimum weighted/positive budget gap over the 34 rows.}
\begin{tabular}{lrrr}
\toprule
relative RMS $\tau$ & common prefix & full prefix & common $>10$\\
\midrule
$1/4$ & 85.3203517096 & 85.3203517096 & $34/34$\\
$1/2$ & 38.2186652435 & 38.2186652435 & $34/34$\\
$3/4$ & 39.2637006403 & 26.0731501545 & $34/34$\\
\bottomrule
\end{tabular}
\end{table}

All 34 source-first weighted ratios are below one and all 34 positive ratios
are above one.  The common weighted budget is above $10^{-5}$ in $102/102$
cases for each source normalizer.  The common-prefix normalization identity
is replayed in 102 cases, and full-prefix tolerance monotonicity in 68
target--row checks.  These counts include the declared source controls and
do not mean that all admissible profiles or shells have been tested.

The finite result is stronger than a mere reuse of the old target atlas: the
17-prime endpoint row, for example, is assigned its label by enumerating all
$2^{16}$ global-sign classes on its own physical Gram matrix.  The result is
still only a stability map.  In particular, the displayed minima do not
constitute a lower bound uniform in shell size, source scale, or profile
cutoff.

\section{Claim firewall and next question}

The exact claims are finite Gram positivity, global-sign reduction,
exhaustive enumeration, budget monotonicity, and normalization cancellation.
The atlas claims only the 34 declared rows, 430 explicit shell targets, 17
profiles, three tolerances, and two target classes.  A uniform native
profile-budget theorem and literal arithmetic $L^2$ estimate remain open;
the fixed-power credit is zero and no full Gate-B or twin-prime result is
claimed.  The natural next project is to test whether the native budget has a
uniform growth law, or to construct the first growing-shell counterexample.

\section*{Reproducibility}

The repository contains the canonical certificate, an independent
source-first replay, exact theorem stress fixtures, and a fail-closed local
Bridge-B checker.  The Session-named Route-A/Route-B evaluator files are not
present in this checkout, so no official evaluator pass is asserted.

\bibliographystyle{plain}
\bibliography{references}
\end{document}
